% ===========================================================================
%  6 — Modelo de dados
% ===========================================================================
\chapter{Modelo de dados}
\label{cap:modelodados}

O sistema mantém dados em três repositórios de naturezas distintas: o banco
relacional (histórico, taxonomia e auditoria), a base vetorial (\textit{chunks}
documentais) e o sistema de arquivos (artefatos de aprendizado de máquina).
Este capítulo especifica os três, com ênfase no relacional, que é onde vive o
modelo de domínio.

\section{Diagrama entidade-relacionamento}

\begin{figure}[H]
\centering
\fitwidth{%
\begin{tikzpicture}[x=1cm, y=1cm]

  \node[entidade, text width=42mm] (ff) at (8.0,0)
    {\etitulo{fault\_families}\nodepart{two}
     \pk{id} \cd{integer PK}\\
     \cd{name varchar(64) UNIQUE}\\
     \cd{is\_fault boolean}};

  \node[entidade, text width=38mm] (lm) at (2.2,-5.0)
    {\etitulo{label\_map}\nodepart{two}
     \pk{raw\_label} \cd{varchar(128) PK}\\
     \fk{family\_id} \cd{integer FK}};

  \node[entidade, text width=48mm] (ev) at (8.0,-5.6)
    {\etitulo{events}\nodepart{two}
     \pk{id} \cd{bigint PK} (sem \textit{autoincrement})\\
     \cd{created\_at timestamptz IDX}\\
     \cd{raw\_fault varchar(128) NULL}\\
     \fk{family\_id} \cd{integer FK NULL IDX}\\
     \cd{z\_rms\_velocity\_mm\_s double}\\
     \cd{x\_rms\_velocity\_mm\_s double}\\
     \cd{temperature\_c double}\\
     \cd{\ldots 15 demais grandezas double}\\
     \cd{rpm double}};

  \node[entidade, text width=38mm] (dc) at (13.8,-5.0)
    {\etitulo{doc\_coverage}\nodepart{two}
     \pk{family\_id} \cd{integer PK/FK}\\
     \pk{document\_id} \cd{integer PK/FK}};

  \node[entidade, text width=42mm] (dg) at (5.2,-11.8)
    {\etitulo{diagnoses}\nodepart{two}
     \pk{id} \cd{bigint PK AUTO}\\
     \cd{event\_id bigint NULL}\\
     \fk{predicted\_family\_id} \cd{int FK}\\
     \cd{probability double}\\
     \cd{neighbors json}\\
     \cd{llm\_response json NULL}\\
     \cd{created\_at timestamptz}};

  \node[entidade, text width=40mm] (dc2) at (13.2,-11.8)
    {\etitulo{documents}\nodepart{two}
     \pk{id} \cd{integer PK AUTO}\\
     \cd{filename varchar(256)}\\
     \cd{title varchar(256)}\\
     \cd{ingested\_at timestamptz}\\
     \cd{status varchar(32)}};

  % --- relacionamentos
  \draw[erel] (ff.west) -- (2.2,-1.6) -- (lm.north)
    node[ercard, pos=0.22, above=1pt]{1}
    node[ercard, pos=0.92, right=2pt]{N};
  \draw[erel] (ff.south) -- (ev.north)
    node[ercard, pos=0.10, right=2pt]{1}
    node[ercard, pos=0.90, right=2pt]{N};
  \draw[erel] (ff.east) -- (13.8,-1.6) -- (dc.north)
    node[ercard, pos=0.22, above=1pt]{1}
    node[ercard, pos=0.92, right=2pt]{N};
  \draw[erel] (dc2.north) -- (dc.south)
    node[ercard, pos=0.10, right=2pt]{1}
    node[ercard, pos=0.90, right=2pt]{N};
  \draw[erel] (ev.south west) -- (5.2,-9.5)
    node[ercard, pos=0.75, right=2pt]{N};
  \node[ercard] at (6.1,-8.2) {1};

  \node[actnote, text width=118mm, anchor=north west] at (1.4,-15.9)
    {\textbf{Relação lógica sem chave estrangeira:} \cd{diagnoses.event\_id} não
     referencia \cd{events} por chave estrangeira, porque um diagnóstico pode ser
     emitido para um evento \emph{ad hoc} --- submetido pela interface e nunca
     persistido no histórico.};
  \draw[actnoteline] (5.2,-15.9) -- (5.2,-13.6);

\end{tikzpicture}}
\caption[Modelo entidade-relacionamento]{Modelo entidade-relacionamento do
  banco corporativo. \cd{fault\_families} é a entidade central: ela ancora a
  taxonomia, os eventos, a cobertura documental e a auditoria de diagnósticos.
  A ausência deliberada de uma chave estrangeira entre \cd{diagnoses} e
  \cd{events} está justificada na nota.}
\label{fig:er}
\end{figure}

\begin{decisao}
A chave primária de \cd{events} é o identificador \textbf{de origem} do evento,
sem geração automática. Essa escolha é o que torna a ingestão MQTT idempotente:
com entrega \textit{at-least-once}, a mesma mensagem pode chegar duas vezes, e
um \cd{merge} pela chave primária resolve o problema na origem, sem tabela de
deduplicação nem janela de tempo (\req{RNF-07}).
\end{decisao}

\section{Dicionário de dados}

\subsection{Tabela \cd{fault\_families}}

Taxonomia canônica de falhas e estados operacionais. É a fonte de verdade da
classificação em todo o sistema --- modelo, API, base documental e interface
usam exatamente estes nomes.

\begin{tabularx}{\linewidth}{@{}L{26mm}L{28mm}L{16mm}X@{}}
\toprule
\headrow \thd{Coluna} & \thd{Tipo} & \thd{Restrição} & \thd{Descrição}\\
\midrule
\cd{id} & \cd{integer} & PK & Identificador da família.\\
\zebra \cd{name} & \cd{varchar(64)} & \textsc{unique} & Nome canônico
  (\textit{snake\_case}), ex.\ \cd{rolamento\_inner}. Usado como chave em todas
  as camadas.\\
\cd{is\_fault} & \cd{boolean} & \textsc{not null} & Distingue falha
  (12 registros) de estado operacional (5 registros). Determina se há prescrição
  a emitir.\\
\bottomrule
\end{tabularx}

\subsection{Tabela \cd{label\_map}}

Registro auditável da canonicalização: preserva a correspondência entre cada
rótulo bruto do conjunto de dados original e a família resultante.

\begin{tabularx}{\linewidth}{@{}L{26mm}L{28mm}L{16mm}X@{}}
\toprule
\headrow \thd{Coluna} & \thd{Tipo} & \thd{Restrição} & \thd{Descrição}\\
\midrule
\cd{raw\_label} & \cd{varchar(128)} & PK & Rótulo bruto normalizado
  (minúsculas, sem espaços nas bordas). São \num{151} valores distintos.\\
\zebra \cd{family\_id} & \cd{integer} & FK & Família canônica atribuída.\\
\bottomrule
\end{tabularx}

\begin{nota}
Esta tabela existe para \emph{auditoria}, não para consulta em tempo de
execução: a canonicalização em \textit{runtime} usa as regras de
\arq{label\_map.yaml}. Manter o resultado materializado permite responder, meses
depois, \enquote{por que este registro foi classificado assim?} sem reexecutar
o ETL.
\end{nota}

\subsection{Tabela \cd{events}}

Histórico de telemetria. É a maior tabela do sistema (\num{166796} registros na
carga inicial) e a base de todas as agregações analíticas.

\begin{tabularx}{\linewidth}{@{}L{40mm}L{22mm}L{15mm}X@{}}
\toprule
\headrow \thd{Coluna} & \thd{Tipo} & \thd{Índice} & \thd{Descrição}\\
\midrule
\endhead
\cd{id} & \cd{bigint} & PK & Identificador de origem do evento. \textbf{Não}
  é gerado pelo banco --- ver decisão acima.\\
\zebra \cd{created\_at} & \cd{timestamptz} & sim & Instante da leitura. Base das
  séries temporais e da frequência diária.\\
\cd{raw\_fault} & \cd{varchar(128)} & --- & Anotação original do operador.
  Identifica também a \emph{campanha de coleta}, fato central na
  Seção~\ref{sec:vazamento}.\\
\zebra \cd{family\_id} & \cd{integer} & sim & Família canônica. Nulo quando o
  evento chega sem anotação (caso normal em produção).\\
\cd{z\_rms\_velocity\_mm\_s} & \cd{double} & --- & Velocidade eficaz no eixo Z
  (\si{\milli\meter\per\second}).\\
\zebra \cd{x\_rms\_velocity\_mm\_s} & \cd{double} & --- & Velocidade eficaz no
  eixo X (\si{\milli\meter\per\second}). Principal indicador de severidade.\\
\cd{temperature\_c} & \cd{double} & --- & Temperatura (\si{\degreeCelsius}).
  \textit{Feature} mais influente do classificador --- e principal suspeita de
  vazamento por sessão.\\
\zebra \cd{z\_peak\_acceleration\_g} & \cd{double} & --- & Aceleração de pico em
  Z ($g$).\\
\cd{x\_peak\_acceleration\_g} & \cd{double} & --- & Aceleração de pico em X
  ($g$).\\
\zebra \cd{z\_peak\_vel\_comp\_freq\_hz} & \cd{double} & --- & Frequência da
  componente de pico de velocidade em Z (\si{\hertz}).\\
\cd{x\_peak\_vel\_comp\_freq\_hz} & \cd{double} & --- & Idem, eixo X. Base da
  \textit{feature} derivada de ordem de rotação.\\
\zebra \cd{z\_rms\_acceleration\_g} & \cd{double} & --- & Aceleração eficaz em Z
  ($g$).\\
\cd{x\_rms\_acceleration\_g} & \cd{double} & --- & Aceleração eficaz em X
  ($g$).\\
\zebra \cd{z\_kurtosis} & \cd{double} & --- & Curtose em Z.\\
\cd{x\_kurtosis} & \cd{double} & --- & Curtose em X.\\
\zebra \cd{z\_crest\_factor} & \cd{double} & --- & Fator de crista em Z.\\
\cd{x\_crest\_factor} & \cd{double} & --- & Fator de crista em X.\\
\zebra \cd{z\_peak\_velocity\_mm\_s} & \cd{double} & --- & Velocidade de pico em
  Z (\si{\milli\meter\per\second}).\\
\cd{x\_peak\_velocity\_mm\_s} & \cd{double} & --- & Velocidade de pico em X
  (\si{\milli\meter\per\second}).\\
\zebra \cd{z\_high\_freq\_rms\_accel\_g} & \cd{double} & --- & Aceleração eficaz
  de alta frequência em Z ($g$).\\
\cd{x\_high\_freq\_rms\_accel\_g} & \cd{double} & --- & Idem, eixo X. Base da
  razão de alta/baixa frequência.\\
\zebra \cd{rpm} & \cd{double} & --- & Rotação do eixo. Variável de condição
  operacional --- estratificar por ela é obrigatório em qualquer leitura de
  severidade.\\
\bottomrule
\end{tabularx}

\subsection{Tabelas \cd{documents} e \cd{doc\_coverage}}

\begin{tabularx}{\linewidth}{@{}L{28mm}L{26mm}L{15mm}X@{}}
\toprule
\headrow \thd{Coluna} & \thd{Tipo} & \thd{Restrição} & \thd{Descrição}\\
\midrule
\multicolumn{4}{@{}l}{\textbf{\cd{documents}}}\\
\cd{id} & \cd{integer} & PK, AUTO & Identificador do documento.\\
\zebra \cd{filename} & \cd{varchar(256)} & \textsc{not null} & Nome do arquivo,
  usado como chave de deduplicação na base vetorial.\\
\cd{title} & \cd{varchar(256)} & \textsc{not null} & Título técnico exibido nas
  citações.\\
\zebra \cd{ingested\_at} & \cd{timestamptz} & \textsc{not null} & Momento da
  indexação.\\
\cd{status} & \cd{varchar(32)} & \textit{default} \cd{active} & Permite retirar
  um procedimento de circulação sem apagar o histórico.\\
\midrule
\multicolumn{4}{@{}l}{\textbf{\cd{doc\_coverage}}}\\
\zebra \cd{family\_id} & \cd{integer} & PK, FK & Família coberta.\\
\cd{document\_id} & \cd{integer} & PK, FK & Documento que a cobre.\\
\bottomrule
\end{tabularx}

\begin{atencao}
\cd{doc\_coverage} é um \textbf{espelho relacional} da cobertura, mantido para
consulta e relatório. A cobertura efetiva --- aquela que o \textit{guardrail}
consulta --- é sempre derivada da presença de \textit{chunks} na base vetorial.
Manter uma única fonte de verdade evita a falha clássica em que a tabela diz
\enquote{coberto} e o índice está vazio.
\end{atencao}

\subsection{Tabela \cd{diagnoses}}

Trilha de auditoria: registra cada diagnóstico emitido, permitindo reconstruir
a posteriori o que o sistema respondeu e com que base.

\begin{tabularx}{\linewidth}{@{}L{32mm}L{22mm}L{15mm}X@{}}
\toprule
\headrow \thd{Coluna} & \thd{Tipo} & \thd{Restrição} & \thd{Descrição}\\
\midrule
\cd{id} & \cd{bigint} & PK, AUTO & Identificador do registro de auditoria.\\
\zebra \cd{event\_id} & \cd{bigint} & \textsc{null} & Evento diagnosticado, se
  houver. Sem FK --- ver nota da Figura~\ref{fig:er}.\\
\cd{predicted\_family\_id} & \cd{integer} & FK, \textsc{null} & Família
  prevista.\\
\zebra \cd{probability} & \cd{double} & \textsc{not null} & Probabilidade
  atribuída pelo classificador.\\
\cd{neighbors} & \cd{json} & --- & Identificadores dos 25 vizinhos que
  sustentaram a resposta.\\
\zebra \cd{llm\_response} & \cd{json} & \textsc{null} & Prescrição gerada e
  citações. Nulo quando a família não é documentada.\\
\cd{created\_at} & \cd{timestamptz} & \textsc{not null} & Momento da emissão.\\
\bottomrule
\end{tabularx}

\section{Esquema \cd{adk} --- memória de conversa}

O \textit{framework} de agentes cria tabelas próprias para persistir sessões de
conversa, e uma delas chama-se \cd{events} --- exatamente o nome da tabela
central do domínio. A colisão é resolvida por isolamento em esquema dedicado.

\begin{tabularx}{\linewidth}{@{}L{40mm}X@{}}
\toprule
\headrow \thd{Aspecto} & \thd{Solução adotada}\\
\midrule
Isolamento & Esquema \cd{adk}, criado na inicialização do agente
  (\cd{CREATE SCHEMA IF NOT EXISTS adk}).\\
\zebra Roteamento & \cd{search\_path} fixado em \cd{adk} nos parâmetros de
  conexão da engine assíncrona, de modo que o \textit{framework} nunca
  \enquote{veja} as tabelas de domínio.\\
Driver & As sessões exigem driver assíncrono; a URL do banco é convertida de
  \cd{psycopg2} para \cd{asyncpg} em tempo de inicialização.\\
\zebra Ciclo de vida & Tabelas criadas pelo próprio \textit{framework}. O
  domínio não depende delas: se o esquema não existir, o agente falha na
  inicialização e o sistema degrada para RAG de um passo.\\
\bottomrule
\end{tabularx}

\section{Base vetorial --- coleção \cd{manuals}}

Uma única coleção, com espaço métrico de cosseno e índice
HNSW~\cite{malkov2020hnsw}. Cada registro é um \textit{chunk} de documento.

\begin{tabularx}{\linewidth}{@{}L{26mm}L{20mm}X@{}}
\toprule
\headrow \thd{Campo} & \thd{Tipo} & \thd{Descrição}\\
\midrule
\cd{id} & texto & Chave composta \cd{\{arquivo\}::\{família\}::\{índice\}}.
  Torna a reindexação idempotente e permite remover um documento por
  filtro.\\
\zebra \textit{embedding} & vetor(768) & Representação densa gerada por
  \cd{gemini-embedding-2} com dimensionalidade truncada.\\
\cd{document} & texto & Conteúdo do \textit{chunk} (aproximadamente
  \num{1200} caracteres).\\
\zebra \cd{metadata.doc} & texto & Nome do arquivo de origem --- vira a citação
  exibida.\\
\cd{metadata.title} & texto & Título técnico do procedimento.\\
\zebra \cd{metadata.section} & texto & Seção corrente no momento da extração
  (ex.\ \cd{4.2 Alinhamento a laser}).\\
\cd{metadata.page} & inteiro & Página de origem.\\
\zebra \cd{metadata.family} & texto & Família canônica coberta. \textbf{É o
  campo que sustenta o guardrail}: a existência de ao menos um registro com
  determinada família \emph{é} a definição de \enquote{documentado}.\\
\bottomrule
\end{tabularx}

\begin{nota}
Um documento que cobre $F$ famílias é indexado $F$ vezes, uma por família. Isso
multiplica o armazenamento, mas permite que o filtro por família seja executado
pelo próprio índice, e não por pós-filtragem --- o que preservaria menos
resultados relevantes para um mesmo \textit{top-k}. Com 268 \textit{chunks} na
base inicial, o custo é irrelevante; em escala de milhares de documentos, a
alternativa seria um campo multivalorado com filtro de inclusão.
\end{nota}

\section{Artefatos de aprendizado de máquina}

\begin{tabularx}{\linewidth}{@{}L{34mm}L{20mm}X@{}}
\toprule
\headrow \thd{Arquivo} & \thd{Tamanho} & \thd{Conteúdo}\\
\midrule
\arq{scaler.joblib} & \SI{1.1}{\kibi\byte} & Médias e desvios das 23
  \textit{features}, ajustados sobre \SI{100}{\percent} dos dados.\\
\zebra \arq{lgbm.joblib} & \SI{38}{\mebi\byte} & Classificador multiclasse de
  400 árvores para 17 classes.\\
\arq{knn.joblib} & \SI{29}{\mebi\byte} & Índice de vizinhos sobre o histórico
  completo padronizado.\\
\zebra \arq{index\_meta.joblib} & \SI{5.8}{\mebi\byte} & Tabela com
  \cd{id}, \cd{created\_at} e \cd{family} de cada linha indexada --- é o que
  traduz posição no índice em evento real.\\
\arq{metrics.json} & \SI{3.9}{\kibi\byte} & Métricas dos três protocolos de
  avaliação, classes, lista de \textit{features} e \textit{timestamp} do
  treino.\\
\bottomrule
\end{tabularx}

Total de \SI{74}{\mebi\byte}, carregados uma única vez por processo em um
\textit{singleton} memorizado. O índice de vizinhos é o componente dominante de
memória em execução e a razão pela qual o serviço de API é o mais pesado da
topologia (\req{RNF-01}).
